\chapter{Doubly Robust Learning}
\label{chap:doubly-robust}

\section*{Outline}

Chapter~\ref{chap:causal-forest} identified conditional treatment effects under
unconfoundedness.  We now return to the average treatment effect (ATE).  Ordinary
least squares need not recover the ATE when treatment is only conditionally random.
We develop outcome-regression and inverse-propensity-weighting representations, then
combine them in the augmented inverse propensity weighted (AIPW) score.  Cross-fitting
makes this score compatible with flexible machine-learning estimates of nuisance
functions.

\section{ATE under unconfoundedness and the failure of OLS}
\label{sec:dr-ols}

Let $W_i$ contain pre-treatment covariates.  We maintain consistency, i.i.d.
sampling, unconfoundedness,
\begin{equation}
  D_i\indep\bigl(Y_i(1),Y_i(0)\bigr)\mid W_i,
  \label{eq:dr-unconfoundedness}
\end{equation}
and overlap,
\begin{equation}
  0<e(W_i)<1\quad\text{almost surely},
  \qquad e(w)=\Pr(D_i=1\mid W_i=w).
  \label{eq:dr-overlap}
\end{equation}
Define the outcome regressions
\begin{equation}
  \mu_d(w)=\E[Y_i\mid D_i=d,W_i=w],\qquad d\in\{0,1\}.
  \label{eq:dr-outcome-regressions}
\end{equation}
By the CATE identification argument,
\begin{equation}
  \tau\equiv\operatorname{ATE}
  =\E[\mu_1(W_i)-\mu_0(W_i)].
  \label{eq:dr-gformula}
\end{equation}
This is often called the outcome-regression formula, the g-formula, or
standardization.

Identification does not imply that an arbitrary linear regression is correctly
specified.  Consider the following data-generating process.  Treatment depends on
$W$, so it is not completely randomized, but the independent variable $V$ makes
treatment random conditional on $W$.  The individual effect is $1+W^2$, hence the
ATE is $1+\E[W^2]=2$.

\begin{rcode}
library(lmtest)
library(sandwich)

set.seed(5280)
n <- 5000
W <- rnorm(n)
V <- rnorm(n)
D <- as.integer(W + V > 0)
epsilon <- rnorm(n)
Y <- 1 + D + W + D * W^2 + sin(W) + epsilon

fit_short <- lm(Y ~ D)
fit_linear <- lm(Y ~ D + W)
coeftest(fit_short, vcov. = vcovHC(fit_short, type = "HC2"))["D", ]
coeftest(fit_linear, vcov. = vcovHC(fit_linear, type = "HC2"))["D", ]
\end{rcode}

Both regressions omit nonlinear terms that are related to treatment.  Their treatment
coefficients are therefore not generally the ATE.  Adding a sufficiently rich and
correctly specified outcome model could solve the problem, but we would like a
method that accommodates flexible estimation and offers protection against one
nuisance model being misspecified.

\section{The propensity score and two representations of the ATE}
\label{sec:dr-identification}

The function $e(W_i)$ in \eqref{eq:dr-overlap} is the \emph{propensity score}.  It
summarizes how observed covariates predict treatment.  Under complete randomization
it is constant; under conditional randomization it may vary with $W_i$.

\subsection{The balancing-score result}

Rosenbaum and Rubin's balancing-score theorem states that
\begin{equation}
  D_i\indep\bigl(Y_i(1),Y_i(0)\bigr)\mid e(W_i)
  \label{eq:dr-ps-unconfoundedness}
\end{equation}
whenever \eqref{eq:dr-unconfoundedness} holds.  To see the key argument, let
$Z_i=(Y_i(1),Y_i(0))$.  Since $D_i$ is binary,
\begin{align*}
  \Pr(D_i=1\mid Z_i,e(W_i))
  &=\E[D_i\mid Z_i,e(W_i)]\\
  &=\E\!\left[\E[D_i\mid Z_i,W_i]\mid Z_i,e(W_i)\right]\\
  &=\E[e(W_i)\mid Z_i,e(W_i)]
   =e(W_i),
\end{align*}
where unconfoundedness gives the third line.  Likewise,
$\Pr(D_i=1\mid e(W_i))=e(W_i)$.  Thus the conditional treatment probability does
not depend on $Z_i$ once the propensity score is held fixed.

The theorem makes a high-dimensional conditioning problem scalar, but it does not
mean that one can ignore $W$ when estimating the propensity score.  Misspecifying
$e(W)$ can still invalidate a propensity-score analysis.

\subsection{Inverse propensity weighting}

The ATE also has the representation
\begin{equation}
  \tau
  =\E\!\left[
      \frac{D_iY_i}{e(W_i)}
      -\frac{(1-D_i)Y_i}{1-e(W_i)}
    \right].
  \label{eq:dr-ipw-id}
\end{equation}
For the treated component, consistency and iterated expectations give
\begin{align*}
  \E\!\left[\frac{D_iY_i}{e(W_i)}\right]
  &=\E\!\left[\frac{D_iY_i(1)}{e(W_i)}\right]\\
  &=\E\!\left[
      \frac{\E[D_iY_i(1)\mid W_i]}{e(W_i)}
    \right]\\
  &=\E\!\left[
      \frac{\E[D_i\mid W_i]\E[Y_i(1)\mid W_i]}{e(W_i)}
    \right]\\
  &=\E[Y_i(1)].
\end{align*}
The third line uses unconfoundedness.  The control component similarly equals
$\E[Y_i(0)]$, proving \eqref{eq:dr-ipw-id}.  Overlap is needed because the
denominators must be defined; near-zero denominators also explain why weak overlap
can make IPW unstable.

We now have two identification formulas: outcome regression in
\eqref{eq:dr-gformula} and IPW in \eqref{eq:dr-ipw-id}.  AIPW combines them.

\section{The AIPW score and double robustness}
\label{sec:dr-aipw}

For generic functions $m_1,m_0$, and $q$ taking values strictly between zero and
one, define
\begin{equation}
  \psi(O_i;m_1,m_0,q)
  =m_1(W_i)-m_0(W_i)
   +\frac{D_i\{Y_i-m_1(W_i)\}}{q(W_i)}
   -\frac{(1-D_i)\{Y_i-m_0(W_i)\}}{1-q(W_i)},
  \label{eq:dr-score}
\end{equation}
where $O_i=(Y_i,D_i,W_i)$.  If the true nuisance functions
$\eta_0=(\mu_1,\mu_0,e)$ were known, the oracle estimator would be
\begin{equation}
  \widehat\tau^*_{\mathrm{AIPW}}
  =\frac1n\sum_{i=1}^n\psi(O_i;\mu_1,\mu_0,e).
  \label{eq:dr-oracle-aipw}
\end{equation}

The augmentation terms have conditional mean zero.  For example,
\[
  \E\!\left[
    \frac{D_i\{Y_i-\mu_1(W_i)\}}{e(W_i)}\Bigm|W_i
  \right]=0.
\]
Hence $\E[\psi(O_i;\eta_0)]=\tau$.  The score can also be rearranged as
\begin{align}
  \psi(O_i;m_1,m_0,q)
  ={}&\frac{D_iY_i}{q(W_i)}
      -\frac{(1-D_i)Y_i}{1-q(W_i)} \notag\\
    &+m_1(W_i)\left\{1-\frac{D_i}{q(W_i)}\right\}
      -m_0(W_i)\left\{1-\frac{1-D_i}{1-q(W_i)}\right\}.
  \label{eq:dr-score-ipw-form}
\end{align}

These two forms reveal \emph{double robustness}:

\begin{itemize}
  \item If $m_d=\mu_d$ for $d=0,1$, then the residual augmentation terms in
  \eqref{eq:dr-score} have conditional mean zero for any admissible $q$.
  Therefore $\E[\psi]=\E[\mu_1(W)-\mu_0(W)]=\tau$.
  \item If $q=e$, then
  \[
    \E\!\left[1-\frac{D}{e(W)}\Bigm|W\right]=0,
    \qquad
    \E\!\left[1-\frac{1-D}{1-e(W)}\Bigm|W\right]=0.
  \]
  Consequently, the last line of \eqref{eq:dr-score-ipw-form} has mean zero for
  any integrable $m_1,m_0$, while its first line identifies the ATE.
\end{itemize}

The conditioning on $W$ in the second argument is essential: an unconditional
mean-zero equality alone would not justify multiplication by an arbitrary function
of $W$.

Let $\widehat\mu_d$ and $\widehat e$ be estimated nuisance functions.  The plug-in
AIPW estimator is
\begin{equation}
  \widehat\tau_{\mathrm{AIPW}}
  =\frac1n\sum_{i=1}^n
  \psi(O_i;\widehat\mu_1,\widehat\mu_0,\widehat e).
  \label{eq:dr-feasible-aipw}
\end{equation}
Subject to regularity conditions, it remains consistent when either both outcome
regressions are consistently estimated or the propensity score is consistently
estimated.  Double robustness protects against one class of nuisance-model
misspecification; it does not protect against violations of unconfoundedness,
overlap, consistency, or sampling assumptions.

\section{Cross-fitting and double machine learning}
\label{sec:dr-dml}

Flexible machine-learning estimates can overfit.  Evaluating a nuisance function on
the same observations used to train it creates dependence that complicates the
first-order expansion of \eqref{eq:dr-feasible-aipw}.  \emph{Cross-fitting} removes
this own-observation link.

Partition the sample into $K$ folds $I_1,\ldots,I_K$.  For each fold $k$:

\begin{enumerate}
  \item estimate $\mu_1$, $\mu_0$, and $e$ using observations outside $I_k$;
  call the resulting functions
  $\widehat\mu_1^{(-k)},\widehat\mu_0^{(-k)},\widehat e^{(-k)}$;
  \item evaluate those functions for observations $i\in I_k$ and compute their
  AIPW scores.
\end{enumerate}

The cross-fitted estimator is
\begin{equation}
  \widehat\tau_{\mathrm{DML}}
  =\frac1n\sum_{k=1}^K\sum_{i\in I_k}
  \psi\!\left(O_i;widehat\mu_1^{(-k)},
  \widehat\mu_0^{(-k)},\widehat e^{(-k)}\right).
  \label{eq:dr-dml}
\end{equation}
Every observation is used for score evaluation once and for nuisance training in
the other folds.

Why can slowly estimated nuisance functions still yield root-$n$ inference for
$\tau$?  The AIPW score is orthogonal: its population expectation is locally
insensitive, to first order, to small perturbations of the nuisance functions.  The
leading remainder is governed by a \emph{product} of nuisance errors.  A typical
sufficient condition is
\begin{equation}
  \|\widehat e-e\|_2
  \left(\|\widehat\mu_1-\mu_1\|_2+
        \|\widehat\mu_0-\mu_0\|_2\right)
  =o_p(n^{-1/2}),
  \label{eq:dr-product-rate}
\end{equation}
together with consistent nuisance estimates, bounded moments, and strong overlap:
$e(W)$ is bounded away from zero and one.

If the propensity and outcome-regression mean-squared errors are respectively of
orders $n^{-a}$ and $n^{-b}$, their $L^2$ errors have orders $n^{-a/2}$ and
$n^{-b/2}$.  Thus \eqref{eq:dr-product-rate} corresponds to $a+b>1$ when these are
exact big-$O$ rates; an equality case needs a stronger little-$o$ statement.  It is
incorrect simply to compare nuisance variances with $1/n$ without taking square
roots and the product remainder into account.

Under appropriate conditions,
\begin{equation}
  \sqrt n(\widehat\tau_{\mathrm{DML}}-\tau)
  \dto \N(0,V_{\mathrm{eff}}),
  \qquad
  V_{\mathrm{eff}}=\Var\{\psi(O_i;\eta_0)\}.
  \label{eq:dr-dml-clt}
\end{equation}
The centered score $\psi(O_i;\eta_0)-\tau$ is the efficient influence function for
the ATE in the nonparametric unconfoundedness model.  Consequently, an estimator
with the expansion in \eqref{eq:dr-dml-clt} attains the semiparametric efficiency
bound among regular estimators in that model.  This is a precise efficiency claim,
not a statement that AIPW is unconditionally ``best'' in every finite sample.

An estimated standard error is obtained from the sample variance of the cross-fitted
scores:
\begin{equation}
  \widehat{\operatorname{se}}(\widehat\tau_{\mathrm{DML}})
  =\sqrt{\frac{1}{n(n-1)}
    \sum_{i=1}^n(\widehat\psi_i-\widehat\tau_{\mathrm{DML}})^2}.
  \label{eq:dr-dml-se}
\end{equation}
The familiar Wald confidence interval then applies when the asymptotic conditions
are credible.

\section{Implementation with generalized random forests}
\label{sec:dr-implementation}

The \texttt{grf} package estimates nuisance components internally and uses
out-of-bag predictions to construct an orthogonalized ATE estimate.  We return to
the simulation from Chapter~\ref{chap:causal-forest}.

\begin{rcode}
library(grf)

set.seed(5280)
n <- 2000
p <- 10
W <- matrix(rnorm(n * p), nrow = n, ncol = p)
D <- rbinom(n, 1, 0.4 + 0.2 * (W[, 1] > 0))
Y <- pmax(W[, 1], 0) * D + W[, 2] + pmin(W[, 3], 0) + rnorm(n)

# Use many trees for stable final inference.
forest <- causal_forest(W, Y, D, num.trees = 4000)
average_treatment_effect(
  forest,
  target.sample = "all",
  method = "AIPW"
)
\end{rcode}

Here the true ATE is
$\E[\max\{W_1,0\}]=1/\sqrt{2\pi}\approx0.399$.  Comparing the estimate with this
truth is a useful simulation check.  Four thousand trees are appropriate for a
substantive analysis but unnecessarily slow for a short browser demonstration; the
live companion uses a smaller forest and labels that choice explicitly.

In applied work, also inspect the estimated propensity scores, sensitivity to
forest tuning, and the effective target population.  Extreme propensities can make
AIPW noisy despite its robustness properties.  If overlap is poor, changing the
estimand, trimming by a pre-specified rule, or collecting better comparison data may
be more defensible than relying on a few enormous weights.

\section{Summary}

\begin{itemize}
  \item Under unconfoundedness and overlap, the ATE has both an outcome-regression
  representation and an inverse-propensity-weighted representation.
  \item OLS with an arbitrary linear specification need not estimate the ATE when
  treatment depends on covariates.
  \item AIPW combines the two representations and is consistent if either the two
  outcome regressions or the propensity score is consistently estimated, subject
  to regularity conditions.
  \item Cross-fitting permits flexible nuisance learners while avoiding
  own-observation overfitting in the score.
  \item Orthogonality makes the leading bias depend on the product of nuisance
  errors.  With strong overlap and suitable product rates, AIPW is root-$n$ normal
  and semiparametrically efficient.
  \item These statistical protections do not repair a failure of the identifying
  causal assumptions.
\end{itemize}
